\documentclass[11pt]{article}
\usepackage[margin=1in]{geometry}
\usepackage{booktabs}
\usepackage{amsmath}
\usepackage{hyperref}
\usepackage{xcolor}
\usepackage{microtype}

\title{Reasoning Scaffold or Statistical Artifact?\\
       Dissecting Chain-of-Thought Gains Across Model Families}
\author{Anonymous}
\date{}

\begin{document}
\maketitle

\begin{abstract}
Chain-of-thought (CoT) prompting reliably improves accuracy on multi-step reasoning benchmarks,
but the mechanism remains contested: does the gain come from genuine semantic reasoning over
intermediate steps, or from statistical conditioning on CoT-shaped tokens regardless of their
truth value? We run a controlled experiment with five prompt conditions and evaluate two model
families on 200 GSM8K problems. Correct traces outperform semantically broken traces by 48--56
percentage points, a gap that survives answer-copying controls, establishing semantic content
as the primary driver. Statistical conditioning is real but secondary, contributing a consistent
$+21$--$24$\,pp lift from broken traces over no-CoT. When given unconstrained output space,
GPT-4o-mini's self-generated CoT achieves 93\%, within 4 points of the verified-trace ceiling.
Trace engagement depth is model-dependent: GPT-4o-mini treats extraneous-premise traces
identically to correct ones, while Grok-4.20 incurs a 5\,pp penalty from the same traces. A
reasoning-native model achieves 75\% on the direct condition with no external trace, reducing
the conditioning-vs-content question to a hierarchy of bottlenecks: weight recall, architectural
reasoning capacity, output budget, and trace quality.
\end{abstract}

%-----------------------------------------------------------------------
\section{Introduction}
%-----------------------------------------------------------------------

Chain-of-thought (CoT) prompting~\citep{wei2022chain} has become a standard technique for
improving accuracy on reasoning benchmarks. The dominant explanation is that explicit
intermediate steps enable the model to work through a problem, accessing computation
unavailable in a single forward pass. This account predicts that step quality should matter:
correct steps should outperform incorrect ones, and any CoT-formatted text should outperform
none.

An alternative account holds that gains arise from \emph{statistical conditioning}: the
presence of reasoning-formatted tokens shifts output distributions toward answer-like
responses, regardless of logical validity. Under this view, semantically broken traces,
correct in format but wrong in intermediate arithmetic, should help nearly as much as
correct ones.

These accounts make different quantitative predictions. If conditioning dominates, CoT
``works'' for reasons unrelated to interpretability or reliability. If semantic content
dominates, CoT traces provide a faithful window into model reasoning, and trace quality
becomes a meaningful signal.

We run a controlled experiment across five conditions to test these predictions, including a
\textsc{permissive-direct} condition where the model generates its own CoT before answering.
The design draws on a typology from the sociology of communication~\citep{latour1984irreductions,attali1984illusions}:
the distinction between \emph{scribe} intermediaries, which preserve information, and
\emph{usurer} intermediaries, which insert extraneous content. This maps directly onto the
space of CoT traces we construct.

%-----------------------------------------------------------------------
\section{Experimental Design}
%-----------------------------------------------------------------------

\subsection{Dataset}

We sample 200 problems uniformly from the GSM8K test set~\citep{cobbe2021gsm8k}, a benchmark
of grade-school math word problems with integer answers and human-written solution traces.

\subsection{Conditions}

\begin{description}
\item[\textsc{direct}] No CoT. The model produces only a final numerical answer.
\item[\textsc{broken}] A solution trace with correct format and token length but deliberately
    wrong intermediate arithmetic. The final answer phrase is formatted correctly but derives
    from erroneous steps.
\item[\textsc{scribe}] A locally valid, information-preserving solution trace: correct
    arithmetic, correct final answer.
\item[\textsc{usurer}] A correct solution trace augmented with extraneous premises or a
    spurious optimization framing before arriving at the correct answer.
\item[\textsc{permissive-direct}] No trace provided. The model is permitted to reason freely
    before answering (500-token budget). This measures whether the model's self-generated
    reasoning reproduces the conditioning effect.
\end{description}

Traces for \textsc{broken}, \textsc{scribe}, and \textsc{usurer} are generated by
\texttt{grok-3-mini} and verified to have matched token lengths.
A \emph{stripped} variant (final answer line removed) rules out answer-copying: accuracy
changed by $< 2$\,pp in both \textsc{scribe} and \textsc{usurer}.

\subsection{Models}

We evaluate GPT-4o-mini (OpenAI Batch API) on all five conditions; Grok-4.20-non-reasoning
(xAI API) on \textsc{direct}, \textsc{broken}, \textsc{scribe}, and \textsc{usurer}; and
Grok-4.20-reasoning (xAI API) on \textsc{direct} only, to test whether built-in
chain-of-thought generation substitutes for external conditioning.

Answers are extracted via a rule-based extractor that prioritizes explicit answer-marker
phrases and falls back to the last number in the response.

%-----------------------------------------------------------------------
\section{Results}
%-----------------------------------------------------------------------

\begin{table}[t]
\centering
\caption{Accuracy (\%) on 200 GSM8K problems by condition and model.
         Grok-4.20-R = reasoning variant (\textsc{direct} only; internal CoT not accessible).}
\label{tab:main}
\begin{tabular}{lccc}
\toprule
Condition & GPT-4o-mini & Grok-4.20-NR & Grok-4.20-R \\
\midrule
\textsc{direct}            & 28 & 12 & \textbf{75} \\
\textsc{broken}            & 49 & 36 & -- \\
\textsc{permissive-direct} & 93 & -- & -- \\
\textsc{usurer}            & 97 & 87 & -- \\
\textsc{scribe}            & 97 & 92 & -- \\
\midrule
$\Delta$\,broken$-$direct  & $+21$ & $+24$ & -- \\
$\Delta$\,scribe$-$broken  & $+48$ & $+56$ & -- \\
\bottomrule
\end{tabular}
\end{table}

Table~\ref{tab:main} presents all results.

\paragraph{Statistical conditioning is real but secondary.}
The broken$-$direct gap is $+21$\,pp (GPT-4o-mini) and $+24$\,pp (Grok-4.20), replicating
across model families. Semantically broken traces, carrying no valid reasoning, improve
accuracy by roughly one-fifth of the total possible gain. This is the pure conditioning signal.

\paragraph{Semantic content is the primary driver.}
The scribe$-$broken gap ($+48$\,pp and $+56$\,pp) is more than twice the conditioning gain
in both cases. This gap is not explained by answer-copying (stripped control: $< 2$\,pp).
Both models genuinely use the semantic content of correct intermediate steps.

\paragraph{Self-generated CoT nearly matches a verified external trace.}
\textsc{permissive-direct} for GPT-4o-mini (93\%) falls within 4\,pp of \textsc{scribe}
(97\%). When allowed to reason freely, the model reaches nearly the same accuracy as when
given a verified correct scaffold. The residual 4\,pp gap is the marginal value of an
external correct trace over unconstrained self-reasoning.

\paragraph{Trace engagement depth varies across architectures.}
GPT-4o-mini treats \textsc{usurer} and \textsc{scribe} identically (both 97\%), suggesting
it uses the trace primarily as a statistical scaffold without critically parsing semantics.
Grok-4.20 shows a 5\,pp penalty for \textsc{usurer} vs.\ \textsc{scribe} (87\% vs.\ 92\%),
indicating deeper engagement with trace content, enough to be misled by extraneous premises.
This difference is not predicted by either pure account.

\paragraph{A reasoning model collapses the direct-condition gap.}
Grok-4.20-R on \textsc{direct} (75\%) far exceeds Grok-4.20-NR on \textsc{direct} (12\%)
and approaches GPT-4o-mini on \textsc{permissive-direct} (93\%). The reasoning model achieves
this with no external trace by generating internal chain-of-thought. Compared to the
non-reasoning variant's \textsc{scribe} ceiling of 92\%, the 17\,pp gap is the marginal
value of a verified correct trace over the model's own unconstrained reasoning.

%-----------------------------------------------------------------------
\section{Discussion}
%-----------------------------------------------------------------------

\subsection{A Decomposition of CoT Gains}

The results support a decomposition of the total CoT gain:
\begin{align}
\underbrace{\text{scribe} - \text{direct}}_{\text{total CoT gain}} =
\underbrace{(\text{broken} - \text{direct})}_{\text{statistical conditioning}}
+ \underbrace{(\text{scribe} - \text{broken})}_{\text{semantic content value}}.
\end{align}
For GPT-4o-mini: $69 = 21 + 48$. For Grok-4.20: $80 = 24 + 56$. Semantic content accounts
for 70\% of total gain in both cases.

\subsection{The Permissive-Direct and Reasoning-Direct Conditions}

\textsc{permissive-direct} (93\%) shows that most of the accuracy gap between \textsc{direct}
(28\%) and \textsc{scribe} (97\%) is achievable without any external trace, simply by
expanding the output budget. The reasoning model (\textsc{direct}, 75\%) adds a further
data point: a model with sufficient architectural capacity can generate effective multi-step
computation internally, even without an explicit prompt inviting reasoning.

The hierarchy direct-NR (12\%) $<$ direct-R (75\%) $<$ permissive-direct (93\%)
$\leq$ scribe (97\%) decomposes the CoT benefit into four levels. The 12\% baseline reflects
weight recall alone. The jump to 75\% reflects architectural reasoning capacity. The jump to
93\% reflects an explicit output budget. The residual gap to 97\% is the value of a verified
external trace over self-generated reasoning.

\subsection{Reconciling with Prior Work}

Prior work has produced mixed evidence on CoT mechanism.
\citet{wang2023towards} find that step validity matters substantially;
\citet{madaan2022text} find that perturbations have limited effect.
The present results reconcile both: conditioning provides a floor ($+21$--$24$\,pp),
semantic content provides the ceiling. The apparent contradiction between prior results
reflects differences in perturbation severity and whether the final answer phrase was
controlled.

\subsection{Implications for CoT Reliability}

Since semantic content is the primary driver, trace quality predicts accuracy. The
Grok-4.20 usurer finding adds nuance: deeper trace engagement increases sensitivity to
correct reasoning but also vulnerability to misleading content. Whether a model reads traces
critically is an architectural property, not a prompting variable, and matters for deployment
in settings where trace quality cannot be guaranteed.

%-----------------------------------------------------------------------
\section{Related Work}
%-----------------------------------------------------------------------

CoT prompting was introduced by~\citet{wei2022chain} and shown to emerge with scale.
Zero-shot CoT~\citep{kojima2022large} demonstrated that even minimal prompts elicit reasoning.
\citet{lanham2023measuring} studied CoT faithfulness; \citet{turpin2023language} showed CoT
can diverge from model computation.
The scribe/usurer typology is adapted from~\citet{latour1984irreductions}; the
conditioning-illusion framing from~\citet{attali1984illusions}.

%-----------------------------------------------------------------------
\section{Conclusion}
%-----------------------------------------------------------------------

CoT gains on GSM8K arise primarily from semantic reasoning capacity. Statistical conditioning
contributes a real but secondary floor of $+21$--$24$\,pp, consistent across model families.
Self-generated CoT reaches 93\%, within 4 points of a verified external trace. A reasoning
model achieves 75\% on the direct condition with no trace at all, collapsing the conditioning
gap and reframing CoT benefit as a hierarchy of bottlenecks rather than a single mechanism.
Trace engagement depth varies across architectures in ways that have direct implications for
when CoT traces can be trusted.

\paragraph{Limitations.} Grok-4.20 \textsc{permissive-direct} was not evaluated; the
comparison between direct-R and permissive-direct within the same model family remains
future work.

\bibliographystyle{plainnat}
\bibliography{refs}

\end{document}
